\chapter{\nt{GLUT} і \nt{GABA} разом}
\label{sec:gaba}

\section{Правила гри}

Мережа з самих лише \nt{GLUT}-нейронів не вміє обирати, скидати пам'ять і втримувати активність від лавини, і все через монотонність (твердження~\ref{prop:monotone}). Додаймо до неї другий за чисельністю тип нейронів --- \nt{GABA}. Правила ті самі: лише те, що буває в живому організмі, і жодного годинника.

Модель нейрона та сама, що~\eqref{eq:glutneuron}, але імпульс від \nt{GABA}-нейрона не піднімає напругу отримувача, а опускає. Ваги тепер бувають двох знаків, і знак задає відправник, правило~\eqref{eq:dalesign}.

Кілька параметрів схеми. \nt{GABA}-нейронів у корі від п'ятої частини до чверті~\cite{Hendry1987}. Виходи в них короткі: вони гальмують сусідів, а не далекі області. Гальмування приходить через мілісекунду й триває кілька мілісекунд. Без входу \nt{GABA}-нейрон мовчить: щоб загальмувати когось, він сам мусить отримати збудження, зазвичай від \nt{GLUT}-нейронів тієї самої мережі.

Тому від'ємний зв'язок від \nt{GLUT}-нейрона $A$ до нейрона $B$ доводиться робити через посередника: $A$ збуджує \nt{GABA}-нейрон, а той гальмує $B$. Зміна знака коштує одного нейрона й однієї затримки, близько мілісекунди.

Для гальмування важливо не лише, \emph{від кого} йде зв'язок, а й \emph{куди} він сідає: місце посадки багато в чому визначає, що зв'язок робить~\cite{Markram2004}. \nt{GLUT}-виходи сідають здебільшого на вхідні гілки отримувача, його \emph{дендрити}, і несуть дані: кожен такий вхід --- один доданок суми. Вихід на тіло нейрона діє на суму цілком: так багато швидких \nt{GABA}-нейронів ділять відповідь отримувача й задають, коли йому спрацювати. Вихід на те місце, де народжується імпульс отримувача, --- останнє слово, заборона на весь вихід одразу. А вихід на закінчення чужого проводу змінює ймовірність викиду $p$ однієї вхідної лінії, не чіпаючи решти~\cite{Rudomin1999}; так, зокрема, працюють ворота болю в розділі~\ref{sec:pain}. За межами мозку виходи сідають на м'язові волокна (розділ~\ref{sec:muscles}) і на органи (розділ~\ref{sec:chemoutput}), а деякі не сідають нікуди й викидають медіатор в об'єм тканини або в кров (розділи~\ref{sec:serocomputer} і~\ref{sec:chemoutput}).

\section{НЕ і заборона}

Чи можна тепер зробити НЕ? Майже. Нейрон, що працює при мовчазному вході, мусить звідкись брати збудження. У живому мозку воно є: на кожен нейрон постійно надходить фонова активність від тисяч інших. Нехай фон тримає нейрон $Y$ у роботі, а вхід $x$ через \nt{GABA}-посередника його гальмує. Це НЕ.

Корисніша, однак, \emph{заборона}: «$a$, але не при $b$»:
\begin{empheq}[box=\keybox]{equation}
y \;=\; a \wedge \bar b .
\label{eq:veto}
\end{empheq}
Нейрон $Y$ отримує збудження від $a$ і гальмування від \nt{GABA}-нейрона, якого збуджує $b$. Фон тут не потрібен: збудження дає сам $a$. Із заборон і АБО збирається все. Наприклад, виключне АБО: $x\oplus y=(x\wedge\bar y)\vee(\bar x\wedge y)$ --- два нейрони-заборони, два \nt{GABA}-посередники й один нейрон АБО.

\begin{keyblock}
\begin{proposition}[Повнота \nt{GLUT}+\nt{GABA}]
\label{prop:gabacomplete}
Мережа з \nt{GLUT}- і \nt{GABA}-нейронів може обчислити будь-яку логічну функцію входів, і кожен біт передається одним проводом. Датчики «вмикання---вимикання» з розділу~\ref{sec:glutcomputer} для цього більше не потрібні. Якщо функція має давати одиницю при всіх мовчазних входах, потрібне джерело фонової активності.
\end{proposition}
\end{keyblock}

Доведення: заборона~\eqref{eq:veto} разом з АБО функціонально повна, якщо є стала одиниця, бо НЕ $x$ --- це заборона «одиниця, але не при $x$». А функції, які при всіх мовчазних входах дають нуль, збираються із заборон і АБО й без одиниці.

\section{Напрямок руху}

Заборона в часі, як і І в часі в розділі~\ref{sec:glutcomputer}, дає детектор напрямку руху.

Нехай два сусідні датчики, $A$ і $B$, стоять на відстані $\Delta x$ один від одного. Вихідний нейрон $Y$ отримує збудження від $B$ і гальмування від $A$ через \nt{GABA}-посередника, із затримкою $d$ (рис.~\ref{fig:direction}). Пляма світла, що біжить зі швидкістю $v$ від $A$ до $B$, проходить $A$ у момент $0$, і гальмування приходить до $Y$ у момент $d$; $B$ вона проходить у момент $\Delta x/v$, і тоді ж приходить збудження. Якщо ці моменти збігаються з точністю до вікна~\eqref{eq:window}, гальмування гасить збудження, і $Y$ мовчить. Це відбувається за швидкості близько
\begin{empheq}[box=\keybox]{equation}
v^* \;=\; \frac{\Delta x}{d} .
\label{eq:vstar}
\end{empheq}
За руху від $B$ до $A$ збудження від $B$ приходить у момент $0$, а гальмування від $A$ --- лише в момент $\Delta x/v+d$, коли $Y$ уже спрацював.

\begin{figure}[t]
\centering
\begin{tikzpicture}[>=Stealth,
  nrn/.style={circle,draw,very thick,fill=blue!6,minimum size=0.9cm,inner sep=0pt},
  gab/.style={nrn,fill=red!10},
  inh/.style={-{Bar[width=7pt]},thick,red!70!black}]
\node[nrn] (A) at (0,2) {$A$};
\node[nrn] (B) at (3,2) {$B$};
\node[gab] (G) at (0.9,0.6) {\small \nt{GABA}};
\node[nrn] (Y) at (3,-0.6) {$Y$};
\draw[->,thick] (A) -- (G);
\draw[inh] (G) -- node[below left=-2pt] {\scriptsize затримка $d$} (Y);
\draw[->,thick] (B) -- (Y);
\draw[->,very thick,red!70!black] (Y) -- (4.6,-0.6) node[right,black] {\small вихід};
\draw[->,thick,orange!85!black] (-0.6,2.9) -- (3.6,2.9) node[right,black] {\small мовчить};
\draw[<-,thick,orange!85!black] (-0.6,3.4) -- (3.6,3.4) node[right,black] {\small відповідає};
\foreach \yy/\dd in {2.9/-1,3.4/1}{
  \foreach \k in {1,2,3}{\draw[orange!70!yellow,line width=0.8pt] (1.5+\dd*0.12+\dd*0.13*\k,\yy+0.1-0.1*\k+0.1) -- ++(\dd*0.25,0);}
  \fill[yellow!70!orange,draw=orange!70!black] (1.5,\yy) circle (0.14);}
\node[font=\scriptsize,align=center] at (1.5,3.95) {пляма світла};
\end{tikzpicture}
\caption{Детектор напрямку руху. $B$ збуджує $Y$, $A$ гальмує його через \nt{GABA}-посередника із затримкою $d$. За руху від $A$ до $B$ зі швидкістю близько $\Delta x/d$ гальмування гасить збудження; за зворотного руху воно запізнюється. Жовта пляма з рисками --- світло, що біжить над датчиками.}
\label{fig:direction}
\end{figure}

У сітківці вибірковість до напрямку руху, схоже, створюється саме так: асиметричним гальмуванням із того боку, звідки рух не повинен викликати відповіді~\cite{Barlow1965}.

\section{Віднімання чи ділення}

Досі гальмування в нас віднімало. Насправді воно влаштоване тонше.

Синапс відкриває канал, тобто вмикає провідність. У збуджувального синапсу рівноважний потенціал лежить далеко вище порога, приблизно на $70\mV$ вище спокою: відкритий канал тягне напругу вгору. У гальмівного \nt{GABA}-синапсу канал пропускає хлор, рівноважний потенціал якого лежить біля спокою, і відкритий канал тягне напругу не вниз, а \emph{до спокою}. Якщо відлічувати напругу від спокою, усталена напруга мембрани з провідністю витоку $g_L$, збуджувальною $g_E$ і гальмівною $g_I$ дорівнює
\begin{empheq}[box=\keybox]{equation}
V_\infty \;=\; \frac{g_E\,E_E}{g_L+g_E+g_I} ,
\label{eq:shunt}
\end{empheq}
де $E_E\approx70\mV$ --- рівноважний потенціал збуджувального синапсу. Це закон Ома для подільника напруги: $g_E$ тягне до $E_E$, а $g_L$ і $g_I$ --- до нуля.

$g_I$ стоїть не в чисельнику зі знаком мінус, а в знаменнику: гальмування не \emph{віднімає} від збудження, а \emph{ділить} його (рис.~\ref{fig:shunt}). Таке гальмування називають шунтувальним: відкриті хлорні канали закорочують мембрану. Для мережі це регулювання підсилення. Якщо $g_I$ пропорційна сумарній активності сусідів, кожен нейрон відповідає не на абсолютну силу свого входу, а на її частку в загальній активності. Таке ділення на загальну активність трапляється в багатьох відділах мозку й вважається однією з його стандартних операцій~\cite{Carandini2012}: та сама мережа працює і за слабкого, і за сильного входу.

Застереження: формула~\eqref{eq:shunt} --- про нейрон, який не видає імпульсів. У нейрона, що спрацьовує, напруга весь час тримається біля порога, струм через гальмівну провідність майже сталий, і на \emph{частоту} імпульсів шунт діє здебільшого відніманням~\cite{HoltKoch1997}. Частота ділиться, якщо нейронові одночасно додати врівноважені збуджувальну й гальмівну фонові провідності, як у тремтливому, врівноваженому режимі живої кори~\cite{Chance2002}. Тож ділення мозок отримує не від одного шунта, а від гальмування разом із фоновим шумом~\cite{Carandini2012}.

\begin{figure}[t]
\centering
\begin{tikzpicture}[>=Stealth,x=1.25cm,y=0.05cm]
\draw[->,gray!70!black] (0,0) -- (5.4,0) node[right,black] {\small $g_E/g_L$};
\draw[->,gray!70!black] (0,0) -- (0,68) node[above,black] {\small $V_\infty$, мВ};
\foreach \v in {20,40,60}{\draw[gray!60!black] (-0.05,\v) -- (0.05,\v) node[left=3pt,font=\scriptsize] {\v};}
\foreach \x in {1,2,3,4,5}{\draw[gray!60!black] (\x,-1.5) -- (\x,1.5) node[below=5pt,font=\scriptsize] {\x};}
\draw[very thick,blue!60!black] plot[domain=0:5,samples=60] (\x,{70*\x/(1+\x)});
\draw[very thick,red!70!black] plot[domain=0:5,samples=60] (\x,{70*\x/(3+\x)});
\draw[thick,densely dashed,gray!50!black] plot[domain=0.56:5,samples=60] (\x,{70*\x/(1+\x)-25});
\node[right,font=\small,blue!60!black] at (5.02,58.3) {без гальмування};
\node[right,font=\small,red!70!black] at (5.02,45.5) {ділить: $g_I=2g_L$};
\node[right,font=\small,gray!40!black] at (5.02,32.5) {віднімає $25\mV$};
\end{tikzpicture}
\caption{Шунтувальне гальмування ділить напругу, а не віднімає від неї. Синя крива --- усталена напруга~\eqref{eq:shunt} без гальмування, червона --- з гальмівною провідністю $g_I=2g_L$. Червона --- та сама синя, розтягнута по горизонталі втричі: вхід ніби поділили на $1+g_I/g_L=3$. Штрихова --- гальмування, що віднімає сталі $25\mV$: крива опускається цілком, нахил не змінюється.}
\label{fig:shunt}
\end{figure}

Є й другий наслідок. Стала часу мембрани $\tau_{\text{еф}}=C/(g_L+g_E+g_I)$, і гальмування її вкорочує. А вікно збігу~\eqref{eq:window} пропорційне $\tau$, тож гальмування, що прийшло разом зі збудженням, \emph{звужує вікно}. Схема, у якій вхід збуджує нейрон безпосередньо й водночас, із затримкою в мілісекунду, гальмує його через \nt{GABA}-посередника, --- одна з найпоширеніших у мозку: нейрон встигає відповісти лише на те, що прийшло в першу мілісекунду-другу~\cite{Pouille2001}.

\section{Вибір}

Тепер головне, чого бракувало \nt{GLUT}-мережі: вибір одного з кількох. Візьмімо дві групи \nt{GLUT}-нейронів із входами $I_1$ і $I_2$. Кожна через своїх \nt{GABA}-посередників гальмує іншу із силою $\beta$ (рис.~\ref{fig:wta}). Для частот $r_1,r_2$ отримуємо
\begin{equation}
\tau\,\frac{\dd r_1}{\dd t} = -r_1 + \bigl[I_1-\beta r_2\bigr]_+ ,\qquad
\tau\,\frac{\dd r_2}{\dd t} = -r_2 + \bigl[I_2-\beta r_1\bigr]_+ ,
\label{eq:wta}
\end{equation}
де $[u]_+=\max(u,0)$: частота не буває від'ємною.

\begin{figure}[t]
\centering
\begin{tikzpicture}[>=Stealth,
  nrn/.style={circle,draw,very thick,fill=blue!6,minimum size=0.9cm,inner sep=0pt},
  gab/.style={nrn,fill=red!10,minimum size=0.75cm},
  inh/.style={-{Bar[width=7pt]},thick,red!70!black}]
\node[nrn] (E1) at (0,0) {$r_1$};
\node[nrn] (E2) at (4,0) {$r_2$};
\node[gab] (G1) at (1.4,1.1) {};
\node[gab] (G2) at (2.6,-1.1) {};
\draw[->,thick] (E1) -- (G1); \draw[inh] (G1) -- (E2);
\draw[->,thick] (E2) -- (G2); \draw[inh] (G2) -- (E1);
\draw[->,thick,blue!60!black] (-1.6,0) node[left,black] {$I_1$} -- (E1);
\draw[->,thick,blue!60!black] (5.6,0) node[right,black] {$I_2$} -- (E2);
\end{tikzpicture}
\caption{Вибір із двох. Дві групи \nt{GLUT}-нейронів (сині) гальмують одна одну через \nt{GABA}-посередників (червоні).}
\label{fig:wta}
\end{figure}

Поки працюють обидва нейрони, система~\eqref{eq:wta} лінійна, і її поведінку задають два власні значення: $-(1+\beta)/\tau$ для однакових відхилень обох частот і $-(1-\beta)/\tau$ для протилежних. За слабкого гальмування, $\beta<1$, обидва значення від'ємні: обидва нейрони працюють, гальмування лише приглушує слабшого. За сильного гальмування, $\beta>1$, друге значення стає додатним: будь-яка різниця між $r_1$ і $r_2$ наростає, доки один нейрон не замовкне зовсім.

\begin{keyblock}
\begin{proposition}[Переможець забирає все]
\label{prop:wta}
Якщо взаємне гальмування сильніше за одиницю, $\beta>1$, стан, у якому працюють обидві групи, нестійкий. Мережа приходить у стан, де працює одна група, а інша мовчить. Переможець із частотою $r_1=I_1$ утримує перемогу, поки $I_2<\beta I_1$.
\end{proposition}
\end{keyblock}

Ми перевірили це на моделі. При $\beta=0.5$ і входах $1.0$ та $0.9$ обидві групи працюють, із частотами $0.73$ і $0.53$. При $\beta=2$ і тих самих входах друга група мовчить зовсім. Такий вибір має пам'ять: якщо потім поміняти входи місцями, колишній переможець лишається переможцем, бо $1.0<2\cdot0.9$.

Із сотнею груп усе так само: кожна гальмує решту, виживає одна.

\section{Скидання пам'яті}

Пам'ять із розділу~\ref{sec:glutcomputer} вимикалася лише втомою. Додаймо до неї \nt{GABA}-нейрон, який збуджується сигналом «скидання» і гальмує групу. Гальмування опускає криву $f(wr+I)$ на рис.~\ref{fig:bistable}, верхній перетин зникає, і група згасає --- і лишається в нижньому стані, коли скидання закінчилося. Тепер пам'ять записується одним сигналом і стирається іншим, як у тригера зі входами встановлення й скидання.

Ще красивіше з'єднати дві такі групи взаємним гальмуванням, як на рис.~\ref{fig:wta}, і дати кожній зворотний зв'язок на саму себе. Сигнал на вхід однієї з груп переводить комірку в її стан, і комірка пам'ятає останнє рішення. Це прямий аналог засувки з двох перехресно з'єднаних елементів АБО-НЕ.

\section{Стійкість і ритм}

Лишається третя біда \nt{GLUT}-мережі --- лавина. Візьмімо групу \nt{GLUT}-нейронів зі зворотним зв'язком на себе $w_{EE}$ і групу \nt{GABA}-нейронів, яких перша збуджує із силою $w_{IE}$ і які гальмують першу із силою $w_{EI}$. Для частот $E$ та $I$ отримуємо
\begin{equation}
\tau_E\frac{\dd E}{\dd t} = -E + w_{EE}E - w_{EI}I + h, \qquad
\tau_I\frac{\dd I}{\dd t} = -I + w_{IE}E ,
\label{eq:ei}
\end{equation}
де $h$ --- зовнішній вхід~\cite{WilsonCowan1972}. Без гальмування при $w_{EE}>1$ активність зростає без меж: кожен імпульс породжує більше ніж один наступний. З гальмуванням усталена частота
\begin{equation}
E^* \;=\; \frac{h}{1-w_{EE}+w_{EI}w_{IE}}
\label{eq:eistar}
\end{equation}
скінченна, якщо знаменник додатний. Але цього замало: рівновага має ще й притягувати. Із лінійного аналізу~\eqref{eq:ei} випливають дві умови.

\begin{keyblock}
\begin{proposition}[Умови стійкості]
\label{prop:ei}
Рівновага~\eqref{eq:eistar} стійка, якщо гальмування
\begin{enumerate}
\item[(i)] \emph{досить сильне}: $w_{EI}\,w_{IE} > w_{EE}-1$;
\item[(ii)] \emph{досить швидке}: $\tau_I < \tau_E/(w_{EE}-1)$.
\end{enumerate}
Якщо порушено (i), активність наростає лавиною. Якщо виконано (i), але порушено (ii), гальмування запізнюється, й активність починає коливатися.
\end{proposition}
\end{keyblock}

Умова~(i) --- це визначник матриці системи~\eqref{eq:ei}, умова~(ii) --- її слід. Зміст другої зрозумілий кожному, хто налаштовував регулятор: запізнілий зворотний зв'язок не гасить відхилення, а розгойдує його.

\begin{figure}[t]
\centering
\begin{tikzpicture}[>=Stealth]
\draw[->,gray!70!black] (0,0) -- (10.9,0) node[below left,black] {\small $t$, мс};
\draw[->,gray!70!black] (0,0) -- (0,3.3) node[left,black] {\small $E$};
\foreach \t in {0,50,100,150}{\draw[gray!70!black] (\t*0.07,0.06) -- (\t*0.07,-0.06) node[below,font=\scriptsize] {\t};}
\draw[very thick,gray!60!black,densely dashed] plot coordinates {(0.000,0.000) (0.035,0.071) (0.070,0.144) (0.105,0.218) (0.140,0.294) (0.175,0.373) (0.210,0.453) (0.245,0.535) (0.280,0.620) (0.315,0.706) (0.350,0.795) (0.385,0.886) (0.420,0.979) (0.455,1.075) (0.490,1.173) (0.525,1.274) (0.560,1.377) (0.595,1.482) (0.630,1.591) (0.665,1.702) (0.700,1.816) (0.735,1.933) (0.770,2.052) (0.805,2.175) (0.840,2.301) (0.875,2.430) (0.910,2.563) (0.945,2.698) (0.980,2.838) (1.015,2.980) (1.050,3.080) (1.085,3.080) (1.120,3.080) (1.155,3.080) (1.190,3.080) (1.225,3.080) (1.260,3.080) (1.295,3.080) (1.330,3.080) (1.365,3.080) (1.400,3.080) (1.435,3.080) (1.470,3.080) (1.505,3.080) (1.540,3.080) (1.575,3.080) (1.610,3.080) (1.645,3.080) (1.680,3.080) (1.715,3.080) (1.750,3.080) (1.785,3.080) (1.820,3.080) (1.855,3.080) (1.890,3.080) (1.925,3.080) (1.960,3.080) (1.995,3.080) (2.030,3.080) (2.065,3.080) (2.100,3.080) (2.135,3.080) (2.170,3.080) (2.205,3.080) (2.240,3.080) (2.275,3.080) (2.310,3.080) (2.345,3.080) (2.380,3.080) (2.415,3.080) (2.450,3.080) (2.485,3.080) (2.520,3.080) (2.555,3.080) (2.590,3.080) (2.625,3.080) (2.660,3.080) (2.695,3.080) (2.730,3.080) (2.765,3.080) (2.800,3.080) (2.835,3.080) (2.870,3.080) (2.905,3.080) (2.940,3.080) (2.975,3.080) (3.010,3.080) (3.045,3.080) (3.080,3.080) (3.115,3.080) (3.150,3.080) (3.185,3.080) (3.220,3.080) (3.255,3.080) (3.290,3.080) (3.325,3.080) (3.360,3.080) (3.395,3.080) (3.430,3.080) (3.465,3.080) (3.500,3.080) (3.535,3.080) (3.570,3.080) (3.605,3.080) (3.640,3.080) (3.675,3.080) (3.710,3.080) (3.745,3.080) (3.780,3.080) (3.815,3.080) (3.850,3.080) (3.885,3.080) (3.920,3.080) (3.955,3.080) (3.990,3.080) (4.025,3.080) (4.060,3.080) (4.095,3.080) (4.130,3.080) (4.165,3.080) (4.200,3.080) (4.235,3.080) (4.270,3.080) (4.305,3.080) (4.340,3.080) (4.375,3.080) (4.410,3.080) (4.445,3.080) (4.480,3.080) (4.515,3.080) (4.550,3.080) (4.585,3.080) (4.620,3.080) (4.655,3.080) (4.690,3.080) (4.725,3.080) (4.760,3.080) (4.795,3.080) (4.830,3.080) (4.865,3.080) (4.900,3.080) (4.935,3.080) (4.970,3.080) (5.005,3.080) (5.040,3.080) (5.075,3.080) (5.110,3.080) (5.145,3.080) (5.180,3.080) (5.215,3.080) (5.250,3.080) (5.285,3.080) (5.320,3.080) (5.355,3.080) (5.390,3.080) (5.425,3.080) (5.460,3.080) (5.495,3.080) (5.530,3.080) (5.565,3.080) (5.600,3.080) (5.635,3.080) (5.670,3.080) (5.705,3.080) (5.740,3.080) (5.775,3.080) (5.810,3.080) (5.845,3.080) (5.880,3.080) (5.915,3.080) (5.950,3.080) (5.985,3.080) (6.020,3.080) (6.055,3.080) (6.090,3.080) (6.125,3.080) (6.160,3.080) (6.195,3.080) (6.230,3.080) (6.265,3.080) (6.300,3.080) (6.335,3.080) (6.370,3.080) (6.405,3.080) (6.440,3.080) (6.475,3.080) (6.510,3.080) (6.545,3.080) (6.580,3.080) (6.615,3.080) (6.650,3.080) (6.685,3.080) (6.720,3.080) (6.755,3.080) (6.790,3.080) (6.825,3.080) (6.860,3.080) (6.895,3.080) (6.930,3.080) (6.965,3.080) (7.000,3.080) (7.035,3.080) (7.070,3.080) (7.105,3.080) (7.140,3.080) (7.175,3.080) (7.210,3.080) (7.245,3.080) (7.280,3.080) (7.315,3.080) (7.350,3.080) (7.385,3.080) (7.420,3.080) (7.455,3.080) (7.490,3.080) (7.525,3.080) (7.560,3.080) (7.595,3.080) (7.630,3.080) (7.665,3.080) (7.700,3.080) (7.735,3.080) (7.770,3.080) (7.805,3.080) (7.840,3.080) (7.875,3.080) (7.910,3.080) (7.945,3.080) (7.980,3.080) (8.015,3.080) (8.050,3.080) (8.085,3.080) (8.120,3.080) (8.155,3.080) (8.190,3.080) (8.225,3.080) (8.260,3.080) (8.295,3.080) (8.330,3.080) (8.365,3.080) (8.400,3.080) (8.435,3.080) (8.470,3.080) (8.505,3.080) (8.540,3.080) (8.575,3.080) (8.610,3.080) (8.645,3.080) (8.680,3.080) (8.715,3.080) (8.750,3.080) (8.785,3.080) (8.820,3.080) (8.855,3.080) (8.890,3.080) (8.925,3.080) (8.960,3.080) (8.995,3.080) (9.030,3.080) (9.065,3.080) (9.100,3.080) (9.135,3.080) (9.170,3.080) (9.205,3.080) (9.240,3.080) (9.275,3.080) (9.310,3.080) (9.345,3.080) (9.380,3.080) (9.415,3.080) (9.450,3.080) (9.485,3.080) (9.520,3.080) (9.555,3.080) (9.590,3.080) (9.625,3.080) (9.660,3.080) (9.695,3.080) (9.730,3.080) (9.765,3.080) (9.800,3.080) (9.835,3.080) (9.870,3.080) (9.905,3.080) (9.940,3.080) (9.975,3.080) (10.010,3.080) (10.045,3.080) (10.080,3.080) (10.115,3.080) (10.150,3.080) (10.185,3.080) (10.220,3.080) (10.255,3.080) (10.290,3.080) (10.325,3.080) (10.360,3.080) (10.395,3.080) (10.430,3.080) (10.465,3.080) (10.500,3.080)};
\draw[very thick,blue!60!black] plot coordinates {(0.000,0.000) (0.035,0.071) (0.070,0.141) (0.105,0.211) (0.140,0.279) (0.175,0.344) (0.210,0.406) (0.245,0.465) (0.280,0.519) (0.315,0.570) (0.350,0.616) (0.385,0.658) (0.420,0.697) (0.455,0.731) (0.490,0.763) (0.525,0.790) (0.560,0.815) (0.595,0.836) (0.630,0.855) (0.665,0.871) (0.700,0.886) (0.735,0.898) (0.770,0.908) (0.805,0.916) (0.840,0.924) (0.875,0.929) (0.910,0.934) (0.945,0.938) (0.980,0.941) (1.015,0.943) (1.050,0.945) (1.085,0.946) (1.120,0.946) (1.155,0.947) (1.190,0.947) (1.225,0.947) (1.260,0.946) (1.295,0.946) (1.330,0.945) (1.365,0.944) (1.400,0.944) (1.435,0.943) (1.470,0.942) (1.505,0.941) (1.540,0.941) (1.575,0.940) (1.610,0.939) (1.645,0.939) (1.680,0.938) (1.715,0.937) (1.750,0.937) (1.785,0.936) (1.820,0.936) (1.855,0.936) (1.890,0.935) (1.925,0.935) (1.960,0.935) (1.995,0.934) (2.030,0.934) (2.065,0.934) (2.100,0.934) (2.135,0.934) (2.170,0.934) (2.205,0.934) (2.240,0.933) (2.275,0.933) (2.310,0.933) (2.345,0.933) (2.380,0.933) (2.415,0.933) (2.450,0.933) (2.485,0.933) (2.520,0.933) (2.555,0.933) (2.590,0.933) (2.625,0.933) (2.660,0.933) (2.695,0.933) (2.730,0.933) (2.765,0.933) (2.800,0.933) (2.835,0.933) (2.870,0.933) (2.905,0.933) (2.940,0.933) (2.975,0.933) (3.010,0.933) (3.045,0.933) (3.080,0.933) (3.115,0.933) (3.150,0.933) (3.185,0.933) (3.220,0.933) (3.255,0.933) (3.290,0.933) (3.325,0.933) (3.360,0.933) (3.395,0.933) (3.430,0.933) (3.465,0.933) (3.500,0.933) (3.535,0.933) (3.570,0.933) (3.605,0.933) (3.640,0.933) (3.675,0.933) (3.710,0.933) (3.745,0.933) (3.780,0.933) (3.815,0.933) (3.850,0.933) (3.885,0.933) (3.920,0.933) (3.955,0.933) (3.990,0.933) (4.025,0.933) (4.060,0.933) (4.095,0.933) (4.130,0.933) (4.165,0.933) (4.200,0.933) (4.235,0.933) (4.270,0.933) (4.305,0.933) (4.340,0.933) (4.375,0.933) (4.410,0.933) (4.445,0.933) (4.480,0.933) (4.515,0.933) (4.550,0.933) (4.585,0.933) (4.620,0.933) (4.655,0.933) (4.690,0.933) (4.725,0.933) (4.760,0.933) (4.795,0.933) (4.830,0.933) (4.865,0.933) (4.900,0.933) (4.935,0.933) (4.970,0.933) (5.005,0.933) (5.040,0.933) (5.075,0.933) (5.110,0.933) (5.145,0.933) (5.180,0.933) (5.215,0.933) (5.250,0.933) (5.285,0.933) (5.320,0.933) (5.355,0.933) (5.390,0.933) (5.425,0.933) (5.460,0.933) (5.495,0.933) (5.530,0.933) (5.565,0.933) (5.600,0.933) (5.635,0.933) (5.670,0.933) (5.705,0.933) (5.740,0.933) (5.775,0.933) (5.810,0.933) (5.845,0.933) (5.880,0.933) (5.915,0.933) (5.950,0.933) (5.985,0.933) (6.020,0.933) (6.055,0.933) (6.090,0.933) (6.125,0.933) (6.160,0.933) (6.195,0.933) (6.230,0.933) (6.265,0.933) (6.300,0.933) (6.335,0.933) (6.370,0.933) (6.405,0.933) (6.440,0.933) (6.475,0.933) (6.510,0.933) (6.545,0.933) (6.580,0.933) (6.615,0.933) (6.650,0.933) (6.685,0.933) (6.720,0.933) (6.755,0.933) (6.790,0.933) (6.825,0.933) (6.860,0.933) (6.895,0.933) (6.930,0.933) (6.965,0.933) (7.000,0.933) (7.035,0.933) (7.070,0.933) (7.105,0.933) (7.140,0.933) (7.175,0.933) (7.210,0.933) (7.245,0.933) (7.280,0.933) (7.315,0.933) (7.350,0.933) (7.385,0.933) (7.420,0.933) (7.455,0.933) (7.490,0.933) (7.525,0.933) (7.560,0.933) (7.595,0.933) (7.630,0.933) (7.665,0.933) (7.700,0.933) (7.735,0.933) (7.770,0.933) (7.805,0.933) (7.840,0.933) (7.875,0.933) (7.910,0.933) (7.945,0.933) (7.980,0.933) (8.015,0.933) (8.050,0.933) (8.085,0.933) (8.120,0.933) (8.155,0.933) (8.190,0.933) (8.225,0.933) (8.260,0.933) (8.295,0.933) (8.330,0.933) (8.365,0.933) (8.400,0.933) (8.435,0.933) (8.470,0.933) (8.505,0.933) (8.540,0.933) (8.575,0.933) (8.610,0.933) (8.645,0.933) (8.680,0.933) (8.715,0.933) (8.750,0.933) (8.785,0.933) (8.820,0.933) (8.855,0.933) (8.890,0.933) (8.925,0.933) (8.960,0.933) (8.995,0.933) (9.030,0.933) (9.065,0.933) (9.100,0.933) (9.135,0.933) (9.170,0.933) (9.205,0.933) (9.240,0.933) (9.275,0.933) (9.310,0.933) (9.345,0.933) (9.380,0.933) (9.415,0.933) (9.450,0.933) (9.485,0.933) (9.520,0.933) (9.555,0.933) (9.590,0.933) (9.625,0.933) (9.660,0.933) (9.695,0.933) (9.730,0.933) (9.765,0.933) (9.800,0.933) (9.835,0.933) (9.870,0.933) (9.905,0.933) (9.940,0.933) (9.975,0.933) (10.010,0.933) (10.045,0.933) (10.080,0.933) (10.115,0.933) (10.150,0.933) (10.185,0.933) (10.220,0.933) (10.255,0.933) (10.290,0.933) (10.325,0.933) (10.360,0.933) (10.395,0.933) (10.430,0.933) (10.465,0.933) (10.500,0.933)};
\draw[very thick,red!70!black] plot coordinates {(0.000,0.000) (0.035,0.071) (0.070,0.143) (0.105,0.217) (0.140,0.293) (0.175,0.370) (0.210,0.449) (0.245,0.528) (0.280,0.609) (0.315,0.691) (0.350,0.774) (0.385,0.858) (0.420,0.943) (0.455,1.028) (0.490,1.114) (0.525,1.200) (0.560,1.287) (0.595,1.374) (0.630,1.462) (0.665,1.549) (0.700,1.636) (0.735,1.724) (0.770,1.810) (0.805,1.897) (0.840,1.983) (0.875,2.068) (0.910,2.153) (0.945,2.237) (0.980,2.320) (1.015,2.402) (1.050,2.483) (1.085,2.563) (1.120,2.641) (1.155,2.717) (1.190,2.792) (1.225,2.866) (1.260,2.937) (1.295,3.007) (1.330,3.075) (1.365,3.080) (1.400,3.080) (1.435,3.080) (1.470,3.080) (1.505,3.080) (1.540,3.080) (1.575,3.080) (1.610,3.080) (1.645,3.080) (1.680,3.080) (1.715,3.080) (1.750,3.080) (1.785,3.080) (1.820,3.080) (1.855,3.080) (1.890,3.080) (1.925,3.080) (1.960,3.080) (1.995,3.080) (2.030,3.080) (2.065,3.080) (2.100,3.080) (2.135,3.080) (2.170,3.080) (2.205,3.080) (2.240,3.080) (2.275,3.080) (2.310,3.080) (2.345,3.080) (2.380,3.080) (2.415,3.080) (2.450,3.080) (2.485,3.080) (2.520,3.080) (2.555,3.080) (2.590,3.080) (2.625,3.080) (2.660,3.080) (2.695,3.080) (2.730,3.080) (2.765,3.080) (2.800,3.080) (2.835,3.052) (2.870,2.973) (2.905,2.892) (2.940,2.807) (2.975,2.719) (3.010,2.629) (3.045,2.535) (3.080,2.438) (3.115,2.339) (3.150,2.238) (3.185,2.133) (3.220,2.029) (3.255,1.930) (3.290,1.836) (3.325,1.747) (3.360,1.661) (3.395,1.580) (3.430,1.503) (3.465,1.430) (3.500,1.360) (3.535,1.294) (3.570,1.231) (3.605,1.171) (3.640,1.113) (3.675,1.059) (3.710,1.007) (3.745,0.958) (3.780,0.911) (3.815,0.867) (3.850,0.825) (3.885,0.784) (3.920,0.746) (3.955,0.710) (3.990,0.675) (4.025,0.642) (4.060,0.611) (4.095,0.581) (4.130,0.553) (4.165,0.526) (4.200,0.500) (4.235,0.476) (4.270,0.452) (4.305,0.430) (4.340,0.409) (4.375,0.389) (4.410,0.370) (4.445,0.352) (4.480,0.335) (4.515,0.319) (4.550,0.303) (4.585,0.288) (4.620,0.274) (4.655,0.261) (4.690,0.248) (4.725,0.236) (4.760,0.225) (4.795,0.214) (4.830,0.203) (4.865,0.193) (4.900,0.184) (4.935,0.175) (4.970,0.166) (5.005,0.158) (5.040,0.151) (5.075,0.143) (5.110,0.136) (5.145,0.130) (5.180,0.123) (5.215,0.117) (5.250,0.111) (5.285,0.106) (5.320,0.101) (5.355,0.096) (5.390,0.091) (5.425,0.087) (5.460,0.083) (5.495,0.079) (5.530,0.075) (5.565,0.071) (5.600,0.068) (5.635,0.064) (5.670,0.061) (5.705,0.058) (5.740,0.055) (5.775,0.053) (5.810,0.050) (5.845,0.048) (5.880,0.045) (5.915,0.043) (5.950,0.041) (5.985,0.039) (6.020,0.037) (6.055,0.035) (6.090,0.034) (6.125,0.032) (6.160,0.030) (6.195,0.029) (6.230,0.027) (6.265,0.026) (6.300,0.025) (6.335,0.024) (6.370,0.022) (6.405,0.021) (6.440,0.021) (6.475,0.022) (6.510,0.024) (6.545,0.027) (6.580,0.032) (6.615,0.037) (6.650,0.044) (6.685,0.052) (6.720,0.061) (6.755,0.071) (6.790,0.083) (6.825,0.095) (6.860,0.109) (6.895,0.124) (6.930,0.140) (6.965,0.157) (7.000,0.176) (7.035,0.195) (7.070,0.215) (7.105,0.237) (7.140,0.260) (7.175,0.283) (7.210,0.308) (7.245,0.333) (7.280,0.360) (7.315,0.388) (7.350,0.416) (7.385,0.445) (7.420,0.476) (7.455,0.507) (7.490,0.539) (7.525,0.571) (7.560,0.604) (7.595,0.638) (7.630,0.673) (7.665,0.708) (7.700,0.744) (7.735,0.781) (7.770,0.818) (7.805,0.855) (7.840,0.893) (7.875,0.931) (7.910,0.969) (7.945,1.008) (7.980,1.047) (8.015,1.086) (8.050,1.126) (8.085,1.165) (8.120,1.204) (8.155,1.244) (8.190,1.283) (8.225,1.322) (8.260,1.362) (8.295,1.400) (8.330,1.439) (8.365,1.477) (8.400,1.515) (8.435,1.553) (8.470,1.590) (8.505,1.626) (8.540,1.662) (8.575,1.698) (8.610,1.733) (8.645,1.767) (8.680,1.800) (8.715,1.832) (8.750,1.864) (8.785,1.895) (8.820,1.924) (8.855,1.953) (8.890,1.981) (8.925,2.007) (8.960,2.033) (8.995,2.057) (9.030,2.080) (9.065,2.102) (9.100,2.123) (9.135,2.142) (9.170,2.160) (9.205,2.177) (9.240,2.192) (9.275,2.205) (9.310,2.217) (9.345,2.228) (9.380,2.237) (9.415,2.245) (9.450,2.251) (9.485,2.255) (9.520,2.258) (9.555,2.259) (9.590,2.259) (9.625,2.256) (9.660,2.253) (9.695,2.247) (9.730,2.240) (9.765,2.231) (9.800,2.220) (9.835,2.208) (9.870,2.194) (9.905,2.178) (9.940,2.160) (9.975,2.141) (10.010,2.120) (10.045,2.098) (10.080,2.074) (10.115,2.048) (10.150,2.021) (10.185,1.992) (10.220,1.961) (10.255,1.929) (10.290,1.895) (10.325,1.860) (10.360,1.824) (10.395,1.786) (10.430,1.746) (10.465,1.706) (10.500,1.663)};

\node[font=\small,gray!40!black,anchor=west] at (0.9,3.15) {без гальмування: лавина};
\node[font=\small,blue!60!black,anchor=west] at (7.4,1.25) {швидке гальмування};
\node[font=\small,red!70!black,anchor=west] at (7.4,2.55) {повільне гальмування};
\end{tikzpicture}
\caption{Частота \nt{GLUT}-групи зі зворотним зв'язком $w_{EE}=1.5$, $\tau_E=10\ms$, після ввімкнення входу. Без гальмування (штрих) активність зростає без меж. З гальмуванням сили $w_{EI}w_{IE}=2$ і $\tau_I=2\ms$ (синя крива) вона виходить на рівень $E^*=h/(1-1.5+2)=0.67h$. З тим самим гальмуванням, але повільним, $\tau_I=30\ms$ (червона крива), активність коливається.}
\label{fig:ei}
\end{figure}

Вважають, що кора працює саме в такому режимі: її власний зворотний зв'язок досить сильний, щоб без гальмування піти в лавину, і тримає її в робочій точці лише швидке гальмування~\cite{Tsodyks1997isn}.

Цей режим має парадоксальний підпис, за яким його можна впізнати ззовні~\cite{Tsodyks1997isn}. Додаймо до~\eqref{eq:ei} зовнішній вхід $h_I$ безпосередньо \nt{GABA}-групі. Усталені частоти стануть
\begin{equation}
E^* \;=\; \frac{h-w_{EI}\,h_I}{D},\qquad
I^* \;=\; \frac{w_{IE}\,h+(1-w_{EE})\,h_I}{D},\qquad
D \;=\; 1-w_{EE}+w_{EI}w_{IE} .
\label{eq:isn}
\end{equation}
Якщо $w_{EE}>1$, коефіцієнт при $h_I$ у другій формулі від'ємний: підштовхнеш \nt{GABA}-нейрони --- і їхня частота \emph{впаде}. Причина в петлі. Зайве гальмування приглушує \nt{GLUT}-групу, та менше збуджує \nt{GABA}-нейрони, і вони втрачають більше, ніж отримали. За чисел рис.~\ref{fig:ei} ($w_{EE}=1.5$, $w_{EI}w_{IE}=2$, $D=1.5$) кожна одиниця доданого входу знижує $I^*$ на третину одиниці. Цей підпис і знайшли в корі: коли відповідь нейронів зорової кори пригнічується подразником навколо, гальмівний струм, який вони отримують, не зростає, а падає разом зі збуджувальним~\cite{Ozeki2009}. Пізніше той самий підпис знайшли в різних областях і шарах кори~\cite{Sanzeni2020}.

Коливання в разі порушення умови~(ii) у мозку теж трапляються. Петля «\nt{GLUT} збуджує \nt{GABA}, \nt{GABA} гальмує \nt{GLUT}» із затримкою в кілька мілісекунд породжує ритм із періодом порядку $12$--$50\ms$ (частота $20$--$80$\,Гц), і такі швидкі ритми в корі добре відомі~\cite{Cardin2009,Buzsaki2012}. Це не годинник комп'ютера: ритм місцевий і виникає лише тоді, коли мережа активна. Але протягом кількох циклів він ділить час на вікна, у яких нейрони можуть спрацювати.

\section{Підсумок}

\begin{table}[t]
\centering
\small
\begin{tabular}{>{\raggedright}p{3.3cm}>{\raggedright}p{4.4cm}>{\raggedright\arraybackslash}p{4.4cm}}
\toprule
& лише \nt{GLUT} & \nt{GLUT} + \nt{GABA} \\
\midrule
НЕ & лише через датчики «вмикання---вимикання» & заборона $a\wedge\bar b$; НЕ за фонової активності \\
проводів на біт & два & один \\
напрямок руху & немає & заборона із затримкою \\
регулювання підсилення & немає & ділення: шунт разом із фоновим шумом \\
вікно збігу & $\sim\tau$ & звужується гальмуванням \\
вибір одного з багатьох & немає & взаємне гальмування, $\beta>1$ \\
скидання пам'яті & лише втомою & сигналом через \nt{GABA} \\
стійкість & лише втомою & швидкий від'ємний зворотний зв'язок \\
ритм & не потрібен & виникає за повільного гальмування \\
\bottomrule
\end{tabular}
\caption{Що додає \nt{GABA} до мережі з \nt{GLUT}-нейронів.}
\label{tab:gaba}
\end{table}

\nt{GABA} майже не додає мережі нових \emph{обчислень} --- усе це можна було обчислити й із двопровідними датчиками, --- а додає \emph{керування} (таблиця~\ref{tab:gaba}): вибір, скидання, регулювання підсилення й стійкість. Збудження в мозку несе дані, а гальмування вирішує, яким даним пройти, коли й наскільки гучно. Для кожного четвертого-п'ятого нейрона кори це дуже розумний поділ праці.

\section{Задачі}

\begin{exercise}[\lvlA]
\label{ex:03:1}
\emph{Шунт у числах.} $E_E=70\mV$. За~\eqref{eq:shunt} знайдіть $V_\infty$ при $g_E=g_L$ і $4g_L$, без гальмування й із шунтом $g_I=2g_L$. Яке $V_\infty$ при $g_E=4g_L$ дало б віднімання, що знімає при $g_E=g_L$ стільки ж, скільки шунт? У скільки разів шунт звужує вікно збігу?
\end{exercise}

\begin{exercise}[\lvlB]
\label{ex:03:2}
\emph{Виведення умов стійкості.} Знайдіть слід і визначник матриці лінійної системи~\eqref{eq:ei} і виведіть із них обидві умови твердження~\ref{prop:ei}. На межі умови~(ii) рівновага втрачає стійкість коливально. Знайдіть період цих коливань за чисел рис.~\ref{fig:ei}.
\end{exercise}

\begin{exercise}[\lvlB]
\label{ex:03:3}
\emph{Ритм із затримки.} Замінімо повільне гальмування моделі~\eqref{eq:ei} швидким із затримкою $d$: $\tau_E\dot E=-E(t)-GE(t-d)+h$. При $\tau_E=10\ms$ знайдіть найменшу затримку $d_c$, що дає коливання, і їхній період для $G=2$ і $G=5$. Чи потрапляють вони у швидкі ритми кори, $12$--$50\ms$, на відміну від задачі~\ref{ex:03:2}?
\end{exercise}

\begin{exercise}[\lvlB]
\label{ex:03:4}
\emph{Моделювання: вибір із пам'яттю.} Проінтегруйте~\eqref{eq:wta} при $\beta=2$, $\tau=1$. (а)~При $I_1=1$ повільно піднімайте $I_2$ від $0$ до $3$ і опускайте назад: за яких $I_2$ мережа перемикається? (б)~На скільки зростає час рішення з тиші, коли різниця входів $\varepsilon$ зменшується вдесятеро?
\end{exercise}

\begin{exercise}[\lvlB]
\label{ex:03:5}
\emph{Проєктування: детектор на смугу швидкостей.} Детектор рис.~\ref{fig:direction} має мовчати за руху від $A$ до $B$ із часом пробігу $5$--$20\ms$. \nt{GABA}-посередник із затримкою $d_k$ забороняє $Y$ на відрізку $[d_k,\,d_k+5\ms]$ після імпульсу $A$; збудження від $B$ приходить одразу. Скільки потрібно посередників і з якими затримками?
\end{exercise}

\begin{exercise}[\lvlB]
\label{ex:03:6}
\emph{Заборони й фон.} Скільки нейронів потребує парність $x\oplus y\oplus z$ у загальній схемі «\nt{GABA}-посередник на вхід, заборона на кожен набір з $f=1$, спільне АБО»? Побудуйте її з двох нейронів, \nt{GABA} і \nt{GLUT}, що отримують входи безпосередньо, указавши ваги й пороги.
\end{exercise}

\begin{exercise}[\lvlC]
\label{ex:03:7}
\emph{Проєктування: вибір зі ста.} Кожна зі $100$ груп \nt{GLUT} має однаково гальмувати всі інші, але не себе. Лінійні \nt{GABA}-посередники збуджуються групами й гальмують їх; групи можуть збуджувати одна одну, але не себе. Яка найменша кількість посередників? А якщо прямих збуджувальних зв'язків немає зовсім?
\end{exercise}

\begin{exercise}[\lvlC]
\label{ex:03:8}
\emph{Дослідження: коли гальмування парадоксальне.} У моделі рис.~\ref{fig:ei} ($w_{EI}=2$, $w_{IE}=1$) \nt{GLUT}-група отримала передавальну функцію $f(u)=[u]_+^2$, \nt{GABA}-група лишилася лінійною. Вище якого входу $h_c$ добавка на \nt{GABA}-групу знижує її ж частоту? Перевірте моделюванням.
\end{exercise}
